% Part: many-valued-logic
% Chapter: sequent-calculus
% Section: introduction

\documentclass[../../../include/open-logic-section]{subfiles}

\begin{document}

\olfileid{mvl}{seq}{int}

\olsection{Pambuka}

Kalkulus sekuen kanggo logika klasik iku sistem !!{derivation}
kang efisien lan prasaja. Yen logika mawa akeh nilai ditetepake
nganggo matriks kang nilai kabenerane winates, yaiku $V$ winates,
kita bisa menehi kalkulus sekuen kanggo logika kasebut. Gagasan
carane teka saka teges sekuen lan wujud aturan inferensi ing
kasus klasik.

Elinga yen sawijining sekuen
\begin{align*}
    !A_1, \dots, !A_m & \Sequent !B_1, \dots, !B_n
\intertext{bisa ditafsirake minangka !!{formula}}
    (!A_1 \land \cdots \land !A_m) & \lif (!B_1 \lor \cdots \lor
!B_n)
\end{align*}
Kanthi tembung liya, !!^a{valuation}~$\pAssign{v}$ \emph{nyukupi}
sekuen $\Gamma \Sequent \Delta$ yen lan mung yen $\pValue{v}(!A) = \False$
kanggo saweneh $!A \in \Gamma$, utawa $\pValue{v}(!A) = \True$
kanggo saweneh $!A \in \Delta$. Miturut tafsir iki, sekuen wiwitan
$!A \Sequent !A$ tansah kacukupan, amarga $\pValue v(!A) = \True$
utawa $\pValue v(!A) = \False$.

Iki aturan inferensi kanggo kondisional ing~$\Log{LK}$,
tanpa nulis formula sisih $\Gamma$ lan $\Delta$:

\begin{defish}
    \Axiom$ \fCenter !A$
    \Axiom$ !B \fCenter $
    \RightLabel{\LeftR{\lif}}
    \BinaryInf$ !A \lif !B \fCenter $
    \DisplayProof
    \hfill
    \Axiom$ !A \fCenter !B$
    \RightLabel{\RightR{\lif}}
    \UnaryInf$ \fCenter !A \lif !B $
    \DisplayProof
\end{defish}

Yen kita gunakake tafsir semantik sekuen ing ndhuwur,
aturan $\LeftR{\lif}$ bisa diwaca mangkene: yen
$\pValue v(!A) = \True$ lan $\pValue v(!B) = \False$,
mula $\pValue v(!A \lif !B) = \False$. Semono uga,
aturan $\RightR{\lif}$ mratelakake yen $\pValue v(!A)
= \False$ utawa $\pValue v(!B) = \True$, mula
$\pValue v(!A \lif !B) = \True$. Satemene, pratelan
kondisional iki loro-lorone bikondisional. Kanggo aturan
$\LeftR{\land}$ lan $\RightR{\lor}$ ing wujud bakune,
pratelan kondisional kang cocog ora dadi bikondisional.
Nanging ana wujud alternatif aturan iki kang dadi bikondisional:

\begin{defish}
    \Axiom$!A, !B, \Gamma \fCenter \Delta$
    \RightLabel{\LeftR{\land}}
    \UnaryInf$!A \land !B, \Gamma \fCenter \Delta$
    \DisplayProof
    \hfill
    \Axiom$ \Gamma \fCenter \Delta, !A, !B$
    \RightLabel{\RightR{\lor}}
    \UnaryInf$ \Gamma \fCenter \Delta, !A \lor !B$
    \DisplayProof
\end{defish}

Gagasan dhasar iki, yen ditrapake ing logika $n$ nilai,
ngasilake kalkulus sekuen kang nduweni $n$ panggonan tinimbang
loro panggonan, siji kanggo saben nilai kabeneran. Kanggo logika
telung nilai kanthi $V = \{\False, \Undef, \True\}$, sekuen
iku ekspresi $\Gamma \mid \Pi \mid \Delta$. Sekuen iki kacukupan
ing !!a{valuation}~$\pAssign v$ yen lan mung yen
$\pValue{v}(!A) = \False$ kanggo saweneh $!A \in \Gamma$,
utawa $\pValue{v}(!A) = \True$ kanggo saweneh
$!A \in \Delta$, utawa $\pValue{v}(!A) = \Undef$ kanggo
saweneh $!A \in \Pi$. Mula, sekuen wiwitan
$!A \mid !A \mid !A$ tansah kacukupan.

\end{document}
